Pri strojnem učenju obravnavamo \pojem{množico opazovanih objektov} $O$, in
opazujemo lastnosti teh objektov.
\pojem{Lastnost} modeliramo kot preslikavo $V: O \to D_V$, kjer je zaloga
vrednosti $D_V$ lahko bodisi poljubna končna množica (v primeru diskretne
spremenljivke), ali pa podmnožica $\R$, v primeru numerične spremenljivke.
Za diskretne spremenljivke dodatno zahtevamo, da so vrednosti $D_V$ neurejene.
Lastnosti ustrezajo \pojem{spremenljivkam}, ki jih ločimo na dva tipa.
Prve so \pojem{napovedne spremenljivke} $\mb{X} = (X_1, \ldots, X_p)$, ki
predstavljajo podatke, druge pa so \pojem{ciljne spremenljivke}, ki jih želimo
napovedati.
Vsako ciljno spremenljivko lahko obravnavamo posebej, torej si mislimo, da imamo
le eno, $Y$.
Pri nenadzorovanem učenju ciljnih spremenljivk nimamo, tam nas namesto napovedi
zanimajo drugačna vprašanja.

\vprasanje{Opiši podatkovno množico pri strojnem učenju.}

Napovedni model predstavimo s funkcijo $m: D_{X_1} \times D_{X_2} \times \cdots
\times D_{X_p} \to D_Y$.
Funkcija za podane vrednosti $\mb{x}$ izračuna napoved ciljne spremenljivke
$m(\mb{x}) = \hat{y}$.
Napovedne modele ločimo glede na to, kakšno spremenljivko napovedujejo.
\pojem{Regresijski model} napoveduje numerično spremenljivko,
\pojem{klasifikacijski} pa diskretno.

Za mero, kateri napovedni model je boljši od drugega, definiramo \pojem{funkcijo
izgube} $L: D_Y \times D_y \to \zo{0, \infty}$, ki izračuna napako ene napovedi
$\hat{y}$ glede na točno vrednost $y$.
Pri regresijskih modelih običajno uporabljamo kvadratno napako
\[
  L_{\text{SE}}(y, \hat{y}) = {(y - \hat{y})}^2,
\]
pri klasifikacijskih pa
\[
  L_{01}(y, \hat{y}) =
  \begin{cases}
	0 & y = \hat{y}, \\
	1 & \text{sicer}.
  \end{cases}
\]
Napovedna napaka modela je potem povprečna napaka na vseh primerih iz $S$,
\[
  \operatorname{Error}(m, S) = \inv{\abs{S}} \sum_{(\mb{x}, y) \in S} L(y,
  m(\mb{x})).
\]

Pri nadzorovanem učenju razdelimo množico $S$ na učno in testno množico,
$S_{\text{train}}$ in $S_{\text{test}}$, ki sta disjunktni in pokrijeta vse
primere.
Pravimo, da je model \pojem{točen}, če ima majhno napako na učni množici, in
\pojem{splošen}, če ima majhno napako na testni množici.

\vprasanje{Kako definiramo napako napovednega modela? Kaj sta točnost in
  splošnost?}

Optimalen in nepristranski napovedni model bo takšen, ki bo minimiziral
kvadratno napako $E({(Y - m(\mb{X}))}^2)$.
V idealnem svetu bo tak model imel obliko $m^*(\mb{x}) = E(Y \such \mb{X} =
\mb{x})$.
Če predpostavimo, da so vsa opažanja v $S$ enako verjetna, je približek tega
modela funkcija
\[
  m^*(\mb{x}_0) = \inv{\abs{S_0}} \sum_{(\mb{x}, y) \in S_0} y,
\]
kjer je $S_0 = \{(\mb{x}, y) \in S_{\text{train}} \such \mb{x} = \mb{x}_0\}$
množica primerov, kjer je $\mb{x} = \mb{x}_0$.
Problem je v tem, da je v praksi ta množica najverjetneje prazna, ali pa vsebuje
zelo malo primerov.
Namesto tega lahko za $S_0$ vzamemo $k$ najbližjih primerov iz
$S_{\text{train}}$, čemur pravimo \pojem{soseščina} točke $\mb{x}_0$.

\vprasanje{Opiši delovanje metode najbližjih sosedov.}